% ---------------------------------------------------------------------
%  09-tooling.tex -- HPB Markup, the open-source renderer and the
%  prototype in a product management portal. The layout algorithm, the
%  data model and the implementation status are in appendix-a-formal.tex
%  (app:layout, app:impl).
% ---------------------------------------------------------------------
\section{Tooling: HPB Markup and an Open-Source Renderer}
\label{sec:tooling}

% plain listings language for code snippets (HTML, TSX, shell)
\lstdefinelanguage{hpbplain}{}

A notation only spreads if drawing it costs nothing. Every diagram in this
paper was drawn by hand with \texttt{hpb.sty} (\cref{app:package}). That is
fine for a paper. It is not fine for a product line that gets a new branch
every few days: nobody redraws a diagram by hand every week, and a diagram
that is not redrawn is soon wrong.

So HPB comes with tooling. \emph{HPB Markup} is a small text language in
the spirit of Mermaid's \texttt{gitGraph}~\citep{mermaid}: you write a few
lines of text and get an HPB diagram. A set of open-source packages renders
that text as SVG in a plain HTML page, in React, Next.js or Vite, in
Markdown documentation and from the command line. The toolkit is released
under the MIT license; its code, documentation and the canonical
situations of \cref{sec:practice} as test fixtures are published at
\hpbrepo.\footnote{Package names in this section are those of version~1.0
of the toolkit. If a name has to change when the packages are published,
the repository lists the published names.} At the end of this section I
describe the prototype I built inside the product management portal I
work with, which draws the same diagrams directly from the trackers.

% =====================================================================
\subsection{Why a text language}
\label{sec:tooling-why}

Product documentation increasingly lives as text: Markdown pages, README
files, wikis, pull requests. Text has three properties that pictures do not.
It can be \emph{diffed}, so a reviewer sees that \texttt{X-212} was moved
from \texttt{v2.1} to \texttt{v2.2} instead of comparing two images. It can
be \emph{generated}, so a script can turn an export from any tracker into an
HPB diagram. And it can be \emph{checked}, so the rules of the method are
applied every time the diagram is drawn, not only when someone remembers
them.

The last point matters most to me. A generic drawing tool will happily
show a branch that never closes, or a release that ships with open
committed work. An HPB Markup document that does either produces a
diagnostic. The language is not only a faster way to draw: it makes
\cref{rule:origin,rule:closure,rule:time,rule:gate} hard to break by
accident.

% =====================================================================
\subsection{Design principles}
\label{sec:tooling-principles}

\begin{itemize}[leftmargin=1.5em,itemsep=2pt,topsep=3pt]
  \item \textbf{Statement order is time order.} Every statement that
        creates an event gets the next column, like the ordinal columns of
        the notation (\cref{sec:notation}). Dates such as \texttt{@2026-03}
        are optional: they only feed the calendar ticks and are checked to
        be non-decreasing.
  \item \textbf{A sensible default fork, an explicit one for issues.}
        Without \texttt{from}, a branch forks from the latest released
        baseline of its line, which is the default of \cref{def:branch} for
        improvements, tasks and lab work. For an issue, \cref{rule:origin}
        requires the baseline in which it was observed: write it with
        \texttt{from v1.4} unless it is the latest release.
  \item \textbf{One terminal per branch.} A branch ends with a merge or a
        discard, or stays open and is drawn up to NOW
        (\cref{rule:closure}).
  \item \textbf{The rules are checked, not assumed.} Merging into a
        baseline that does not succeed the fork baseline, including the
        fork itself, is an error (\cref{rule:time}). Releasing a planned
        baseline that still has open targeted branches is an error
        (\cref{rule:gate}), unless the release is marked \texttt{force}: the
        violation is then recorded and stays visible as a warning badge.
  \item \textbf{Forgiving to write.} Keywords are case-insensitive and have
        aliases (\texttt{bug} for \texttt{issue}, \texttt{feature} for
        \texttt{improvement}); identifiers can be any version or reference
        such as \texttt{v1.0.1} or \texttt{Rev-B}; \texttt{\#} and
        \texttt{\%\%} start comments; indentation is ignored.
  \item \textbf{Clear errors.} Every diagnostic has a stable code, a line,
        a column and a plain message. By default the renderer draws an
        error box instead of throwing, so a broken diagram never breaks the
        page around it.
  \item \textbf{One visual contract.} The renderer draws the glyphs and
        palette of \cref{sec:notation}, the same ones \texttt{hpb.sty}
        draws, so a diagram on a web page and one in a printed report read
        the same way.
\end{itemize}

% =====================================================================
\subsection{A complete example}
\label{sec:tooling-example}

\Cref{fig:markup-example} shows a firmware line written in HPB Markup and
the diagram it produces. It is the same situation as
\cref{fig:software-line}, with different illustrative items. Line by line:

\begin{itemize}[leftmargin=1.5em,itemsep=1pt,topsep=3pt]
  \item \texttt{hpb} gives the diagram a title, and \texttt{line} declares
        the product line \emph{Product A}.
  \item \texttt{release v1.0} creates the first released baseline.
  \item \texttt{improvement I-21} and \texttt{issue X-101} open two
        branches. Neither says where it starts, so both fork from
        \texttt{v1.0}, the latest release. The improvement goes above the
        main line, the issue below it. X-101 has no title: its short life
        leaves no room for a label (layout rule~\ref{lay:labels}).
  \item \texttt{release v1.0.1 merges X-101} publishes a patch that closes
        the defect. The next two branches fork from \texttt{v1.0.1}, now
        the latest release.
  \item \texttt{release v1.1} and \texttt{release v1.2} close the remaining
        work. T-8 runs across two releases on the task lane, below the
        issue lane.
\end{itemize}

\noindent Two layout conventions of the renderer are worth stating. First,
the merge terminals of all branches closed by one
\texttt{release \dots{} merges} statement share one column, just before
the baseline; they sit on different lanes, which is the exception allowed
by layout rule~\ref{lay:ordinal}. Second, NOW is drawn only when the
diagram has something in the future (planned baselines, open branches, or
merges awaiting release); a pure history such as this one has no NOW
marker. The renderer stacks extra lanes per type, outward from the main
line, and packs them first-fit; the preference of
rule~\ref{lay:onelane} for putting still-open branches on the outer lanes
is not implemented yet (the full algorithm is in \cref{app:layout}).

% F-M -- HPB Markup: the canonical example A (a firmware line) as text and
% as the diagram it produces. Illustrative data. Requires hpb.sty v1.1 and
% the listings language "hpb" (defined in main.tex).
% Column positions are those computed by @hpb/core 1.0 for this source
% (x positions divided by the column width): one column per event in
% statement order; the merge terminals of the branches closed by one
% "release ... merges" statement share one column just before the baseline;
% a dated baseline after a shared terminal column gets a small spacer. The
% document is pure history, so the renderer draws no NOW marker.
%   1 v1.0 | 2 I-21 | 3 X-101 | 4 terminal X-101 | 5 v1.0.1 | 6 X-102 |
%   7 T-8 | 8 terminals I-21, X-102 | 9.4 v1.1 | 10.4 I-22 | 11.4 X-103 |
%   12.4 terminals T-8, I-22, X-103 | 14.4 v1.2
\begin{figure}[tbp]
\centering
\begin{minipage}{0.66\linewidth}
\begin{lstlisting}[style=hpb,basicstyle=\ttfamily\footnotesize,numbers=left,numberstyle=\tiny\color{black!45},numbersep=6pt,xleftmargin=16pt]
hpb "Product A: firmware line"
line A "Product A"
release v1.0 @2026-01
improvement I-21 "Faster pairing"
issue X-101
release v1.0.1 merges X-101
issue X-102 "Clock drift"
task T-8 "Migrate logs"
release v1.1 @2026-03 merges I-21, X-102
improvement I-22 "Low-power mode"
issue X-103 "Menu translation"
release v1.2 @2026-05 merges T-8, I-22, X-103
\end{lstlisting}
\end{minipage}

\medskip
\hpbfit{%
\begin{hpbgraph}
  % lanes: improvements above, issues and tasks below the main line
  \hpblane{1}{imp}{Improvement (I)}{0}{15}
  \hpblane{-1}{iss}{Issue (X)}{0}{15}
  \hpblane{-2}{task}{Task (T)}{0}{15}
  % main line: pure history, no NOW marker
  \hpbmainline{0}{15}
  \hpborigin{0}{Product A}
  % released baselines
  \hpbrelease{1}{v1.0}
  \hpbrelease{5}{v1.0.1}
  \hpbrelease{9.4}{v1.1}
  \hpbrelease{14.4}{v1.2}
  % branches (fork col, lane, start col, end col, end, merge col)
  \hpbbranch[label={I-21 Faster pairing}]{imp}{1}{1}{2}{8}{merge}{9.4}
  \hpbbranch[label={X-101}]{iss}{1}{-1}{3}{4}{merge}{5}
  \hpbbranch[forkoff=0.6, label={X-102 Clock drift}]{iss}{5}{-1}{6}{8}{merge}{9.4}
  \hpbbranch[forkoff=0.3, mergeoff=0.3, label={T-8 Migrate logs}]{task}{5}{-2}{7}{12.4}{merge}{14.4}
  \hpbbranch[label={I-22 Low-power mode}]{imp}{9.4}{1}{10.4}{12.4}{merge}{14.4}
  \hpbbranch[mergeoff=0.6, label={X-103 Menu translation}]{iss}{9.4}{-1}{11.4}{12.4}{merge}{14.4}
  % ordinal time axis with the calendar ticks given by the @dates
  \hpbaxis{0}{15}{-2.75}
  \hpbtick{1}{Jan 2026}{-2.75}
  \hpbtick{9.4}{Mar}{-2.75}
  \hpbtick{14.4}{May}{-2.75}
\end{hpbgraph}}
\caption{HPB Markup in practice (illustrative). Top: twelve lines of text
that describe a firmware line. Bottom: the diagram they produce, drawn here
with \texttt{hpb.sty} at the column positions the renderer computes for
this source. Nobody types a coordinate: statement order gives the columns,
the branch type gives the lane, and each \texttt{release \dots{} merges}
closes the listed branches (\cref{rule:origin,rule:closure}). The terminals
closed by one release share a column; the history has no NOW marker.}
\label{fig:markup-example}
\end{figure}


The other canonical situations are equally short. A release gate, a
discard with its rationale, and a hardware revision that forks a new
reference are written like this (the full language is in
\cref{app:markup}):
\begin{lstlisting}[style=hpb]
now
plan v2.1 @2026-12 targets I-50, X-211, X-212
discard I-60 "not viable: BOM cost"
release RevB hardware merges L-5, I-30
line RB "DMS-4000 Rev B" from RevB
\end{lstlisting}

\paragraph{What the Markup checks, and what it does not.}
The Markup checks a structural version of the gate rule: no branch
targeted at a baseline may still be open when that baseline is released.
It does not model verification evidence, so ``verified'' in the sense of
\cref{rule:gate} remains a human check at the gate review. A
\texttt{release \dots{} merges} statement plays the role of publication
(\cref{def:gate}), with the listed branches as its targeted set. A
\texttt{force} release records a deliberate violation: the resulting model
is not well formed under \cref{rule:gate}, and the diagram says so with a
\texttt{W-R4-FORCED} warning and a ``!'' badge.

% =====================================================================
\subsection{Packages}
\label{sec:tooling-packages}

The renderer is a small set of TypeScript packages
(\cref{tab:packages}). The core has no runtime dependencies and produces a
plain SVG string, so it runs the same way in a browser, in Node.js and in
server-side rendering.

\begin{table}[tbp]
  \centering
  \small
  \caption{HPB Markup toolkit, version 1.0 (MIT license).}
  \label{tab:packages}
  \renewcommand{\arraystretch}{1.12}
  \begin{tabularx}{\linewidth}{@{}>{\raggedright\arraybackslash}p{2.6cm}
      >{\raggedright\arraybackslash}X
      >{\raggedright\arraybackslash}p{3.5cm}@{}}
    \toprule
    \textbf{Package} & \textbf{What it does} & \textbf{Typical host} \\
    \midrule
    \texttt{@hpb/core}
      & \texttt{parse}, \texttt{validate}, \texttt{buildModel} (checks
        R1--R4, returns diagnostics), \texttt{layout}, \texttt{renderSVG}
        (pure string, no DOM).
      & Any JavaScript runtime, SSR, build scripts. \\
    \addlinespace
    \texttt{@hpb/web}
      & Mermaid-style auto-render of \texttt{pre.hpb} and
        \texttt{code.language-hpb} blocks; custom element
        \texttt{<hpb-diagram>}; API \texttt{initialize}, \texttt{run},
        \texttt{render}; \texttt{dist/hpb.min.js} exposes
        \texttt{window.HPB} for a \texttt{<script>} tag.
      & Plain HTML pages, wikis, static sites. \\
    \addlinespace
    \texttt{@hpb/react}
      & \texttt{<HpbDiagram code=\dots{} />}, SSR-safe; optional hover
        highlighting, \texttt{onNodeClick} and \texttt{onError}.
      & React, Next.js (App Router), Vite. \\
    \addlinespace
    \texttt{@hpb/cli}
      & \texttt{hpb render} to SVG and \texttt{hpb validate} for
        continuous integration.
      & Terminals, CI pipelines, document builds. \\
    \addlinespace
    \texttt{remark-hpb}
      & remark (unified) and markdown-it plugins that render
        \texttt{hpb} fenced code blocks as inline SVG.
      & Markdown and MDX, documentation sites, static generators. \\
    \bottomrule
  \end{tabularx}
\end{table}

\paragraph{In a plain HTML page.}
Load one script and write the diagram as text. Every matching block is
replaced by its SVG when the page loads.
\begin{lstlisting}[language=hpbplain]
<script src="https://unpkg.com/@hpb/web@1.0.0/dist/hpb.min.js"></script>
<pre class="hpb">
line A "Product A"
release v2.0
issue X-211 "Battery drain"
release v2.0.1 merges X-211
</pre>
\end{lstlisting}

\paragraph{In React, Next.js or Vite.}
The component renders the SVG string during render and never touches
\texttt{window}, so it works in server-side rendering. In the Next.js App
Router, a server component can call \texttt{renderSVG} directly; the
interactive component is a client component.
\begin{lstlisting}[language=hpbplain]
import { HpbDiagram } from "@hpb/react";

<HpbDiagram code={source} theme="default"
            onNodeClick={(n) => openItem(n.id)} />
\end{lstlisting}

\paragraph{From the command line.}
\texttt{hpb validate} returns a non-zero exit code when a document breaks
a rule, which makes it usable as a check in a documentation pipeline.
\begin{lstlisting}[language=hpbplain]
hpb render product-a.hpb -o product-a.svg --theme mono
hpb validate product-a.hpb
\end{lstlisting}

\paragraph{Accessibility and print.}
The SVG carries a title and a description, every node has its own
\texttt{<title>}, and no text is smaller than 10\,px at scale one. The
\texttt{mono} theme renders in grayscale with letters and line styles only,
which is the redundant encoding of \cref{sec:notation} put to work.

% =====================================================================
\subsection{Prototype implementation in a product management portal}
\label{sec:implementation}

A method that needs a new tool, and a migration of all existing work into
it, rarely gets adopted. I did not want HPB to be one more system to keep
in sync, so I built it as a thin layer on top of the trackers in which
improvements, defects, tasks and laboratory evaluations were already
managed, as a module of the product management portal I use daily.

\paragraph{An overlay, not a migration.}
The trackers stay exactly as they are and remain the system of record for
the work itself. HPB adds a few tables that store only what the trackers
do not know: the product lines and their family tree, the baselines, the
attachment of each tracker item to a line as a branch with its target and
terminal, and comments attached to exactly one node
(\cref{app:impl}). Because tracker items are referenced and never copied,
adopting HPB needs no migration, and removing it leaves the trackers
untouched.

\paragraph{Derived mode: a picture on day one.}
The same property gives a read-only \emph{derived mode} for products
whose overlay is still empty. Baselines are rebuilt from existing release
records, branches from the tracker items linked to the product, and
terminals from tracker statuses. A completed item with no link to a
release is drawn as closed without a baseline: a violation of
\cref{rule:closure} that the diagram shows instead of hiding. On the first
day, derived mode gives a retrospective picture of the product and a list
of the gaps the method then closes.

\paragraph{The diagram as the place where decisions happen.}
In the portal the diagram is not just a view. One action on a line
creates a typed item in its tracker and attaches it as a branch, so it has
a type and an origin from the first second
(\cref{rule:origin,rule:typed}). Targeting offers only planned baselines;
a discard asks for a rationale; publishing a planned baseline lists
everything targeted at it and merges the set in one operation; a hardware
baseline that changes the reference records the new line it creates
(\cref{rule:lineage}); and every line, baseline and branch has its own
comment thread and a link back to the item in its tracker
(\cref{rule:discussion}). Several rule checks are designed but not yet
enforced in the prototype; \cref{app:impl} lists them honestly. Until they
are, the violations are at least visible in the diagram, and HPB Markup
already enforces R1--R4 for anyone who draws with it.
